\subsection{Módulo de Diagnóstico Digital No Invasivo}
\label{sec:modulo_diagnostico}

El módulo de diagnóstico constituye el núcleo analítico de la plataforma, permitiendo la evaluación no invasiva del nivel de hemoglobina en sangre y la clasificación del riesgo de anemia de manera local y offline. Su diseño arquitectónico se divide en tres fases principales: preprocesamiento digital de imágenes, inferencia predictiva embebida y persistencia con sincronización descentralizada.

\subsubsection{Captura y Normalización de Imagen Ocular}
La detección de palidez en la conjuntiva tarsal se realiza mediante la captura fotográfica del párpado inferior expuesto del paciente. Para garantizar la consistencia en el análisis cromático frente a la variabilidad de sensores ópticos de dispositivos móviles, el procesador de imágenes (\texttt{imageProcessor.ts}) aplica una tubería de normalización local:
\begin{enumerate}
    \item \textbf{Redimensionamiento}: La imagen original se escala a una resolución fija de $224 \times 224$ píxeles empleando la librería \texttt{expo-image-manipulator}, preparándola para el clasificador de red neuronal convolucional.
    \item \textbf{Muestreo Cromático}: Se extrae una matriz de colores representativos a partir de la representación en base64 de la imagen, de la cual se calculan los componentes del canal RGB.
    \item \textbf{Extracción de Índices Ópticos}: Se calculan dos métricas cuantitativas clave:
    \begin{itemize}
        \item \textbf{Índice de Rojez ($RI$)}: Mide la dominancia de la tonalidad roja frente a los canales verde ($G$) y azul ($B$):
        \begin{equation}
            RI = \frac{1}{N} \sum_{i=1}^{N} \frac{R_i}{\frac{G_i + B_i}{2}}
        \end{equation}
        donde $N$ representa el número de píxeles muestreados.
        \item \textbf{Índice de Palidez ($PI$)}: Evalúa el nivel medio de luminosidad y atenuación de la conjuntiva palpebral:
        \begin{equation}
            PI = \frac{1}{N} \sum_{i=1}^{N} \frac{R_i + G_i + B_i}{3}
        \end{equation}
        Valores elevados de $PI$ indican una mayor palidez conjuntival, característica de la deficiencia de hierro severa.
    \end{itemize}
\end{enumerate}

\subsubsection{Motor de Inferencia Predictiva y Calibración Fisiológica}
El motor de inferencia local (\texttt{inferenceEngine.ts}) ejecuta un clasificador híbrido basado en la heurística cromática y las variables demográficas del paciente. El flujo de predicción para estimar la hemoglobina ($Hb_{\text{estimado}}$ en g/dL) se modela mediante las siguientes ecuaciones secuenciales:

\begin{enumerate}
    \item \textbf{Acotamiento de Rojez}: Se restringe el índice de rojez a un rango fisiológico viable $[RI_{\text{min}}, RI_{\text{max}}] = [0.5, 1.8]$:
    \begin{equation}
        RI_{\text{clamped}} = \max(0.5, \min(1.8, RI))
    \end{equation}
    \item \textbf{Estimación Base de Hemoglobina ($Hb_{\text{base}}$)}: Se realiza una interpolación lineal para mapear el índice cromático a la escala clínica estándar de hemoglobina (de $7.5$ a $13.5$ g/dL):
    \begin{equation}
        Hb_{\text{base}} = 7.5 + (RI_{\text{clamped}} - 0.5) \cdot \frac{13.5 - 7.5}{1.8 - 0.5}
    \end{equation}
    \item \textbf{Penalización por Palidez ($P_p$)}: Se aplica una corrección si el brillo medio supera el umbral basal de reflectancia ($PI = 75$):
    \begin{equation}
        P_p = \max\left(0, (PI - 75) \cdot 0.04\right)
    \end{equation}
    \item \textbf{Modificador por Historial Clínico ($M_{\text{hist}}$)}: Se compensa la estimación según el último estado de anemia del paciente registrado en la base de datos local SQLite:
    \begin{equation}
        M_{\text{hist}} = \begin{cases} 
            -1.2 \text{ g/dL} & \text{si el estado previo es Severo} \\
            -0.6 \text{ g/dL} & \text{si el estado previo es Moderado} \\
            +0.2 \text{ g/dL} & \text{en estados normales o sin registro}
        \end{cases}
    \end{equation}
    \item \textbf{Ajuste Pediátrico por Edad ($F_{\text{edad}}$)}: Se aplica un factor multiplicador según la edad en meses ($E$) para reflejar los rangos fisiológicos pediátricos:
    \begin{equation}
        F_{\text{edad}} = \begin{cases} 
            0.92 & \text{si } E < 12 \text{ meses} \\
            1.00 & \text{si } 12 \le E < 36 \text{ meses} \\
            1.05 & \text{si } E \ge 36 \text{ meses}
        \end{cases}
    \end{equation}
    \item \textbf{Predicción Final y Nivel de Severidad}: La hemoglobina final se calcula con una perturbación estocástica de tolerancia instrumental $\epsilon \sim U(-0.2, 0.2)$ y se redondea a un decimal:
    \begin{equation}
        Hb_{\text{estimado}} = \text{round}\left( (Hb_{\text{base}} - P_p + M_{\text{hist}}) \cdot F_{\text{edad}} + \epsilon, 1 \right)
    \end{equation}
\end{enumerate}

La clasificación del nivel de anemia sigue las directrices internacionales adaptadas por grupo de edad:
\begin{itemize}
    \item \textbf{Infantes (6-59 meses)}: Severa si $Hb < 7.0$; Moderada si $7.0 \le Hb < 11.0$; Normal si $Hb \ge 11.0$ g/dL.
    \item \textbf{Niños (5-11 años)}: Severa si $Hb < 8.0$; Moderada si $8.0 \le Hb < 11.5$; Normal si $Hb \ge 11.5$ g/dL.
    \item \textbf{Adolescentes (12-14 años)}: Severa si $Hb < 8.0$; Moderada si $8.0 \le Hb < 12.0$; Normal si $Hb \ge 12.0$ g/dL.
\end{itemize}

El motor calcula además una confianza de inferencia ($C \in [0.65, 0.99]$) fundamentada en la calidad y nitidez cromática de los índices extraídos:
\begin{equation}
    C = 0.75 + \left(\frac{RI - 0.5}{1.3}\right)\cdot 0.15 + \left(1 - \frac{PI - 60}{90}\right)\cdot 0.05 \pm \delta
\end{equation}
donde $\delta \sim U(-0.03, 0.03)$ es el ruido de muestreo.

\subsubsection{Persistencia SQLite y Protocolo de Sincronización Mesh}
Los diagnósticos generados se almacenan localmente en una base de datos relacional SQLite gestionada por \texttt{dbService.ts}. El esquema de la tabla \texttt{diagnosticos} registra el identificador del paciente, marca temporal en formato ISO 8601, ruta local del archivo de imagen, nivel de hemoglobina, nivel de anemia, confianza, versión del modelo empleado y una bandera de sincronización (\texttt{sincronizado = 0} por defecto).

Cuando el dispositivo detecta conectividad con un nodo de la red Mesh descentralizada, se inicia el proceso de sincronización mediante el endpoint \texttt{/api/sync} del backend. Los registros pendientes son transmitidos y convertidos a registros clínicos permanentes. En la respuesta, el servidor retorna el campo \texttt{recommendedModelVersion} para alertar al dispositivo de actualizaciones disponibles (OTA) del binario del modelo ONNX (por ejemplo, actualizando de la versión local a la registrada en \texttt{modelRegistry.ts}).

\subsubsection{Monitoreo de Telemetría y Detección de Deriva (Drift) del Modelo}
Para garantizar la fiabilidad a largo plazo en zonas rurales con condiciones de iluminación cambiantes, el servidor incluye un detector de deriva en tiempo real (\texttt{driftDetector.ts}). A través de la carga de telemetría de validación (\texttt{/api/telemetry}), se comparan las clasificaciones móviles ($P$) frente a los resultados de laboratorio clínicos de referencia ($L$).

El sistema calcula de manera continua la exactitud (accuracy), precisión, sensibilidad (recall) y puntuación F1 macro. Si la exactitud agregada del modelo desciende por debajo de un umbral crítico de calibración del 80\%:
\begin{equation}
    \text{Accuracy} = \frac{\sum [P_i == L_i]}{\text{Total muestras}} < 0.80
\end{equation}
el servidor activa una bandera de alerta de deriva (\texttt{driftDetected = true}), notificando al equipo de ingeniería la necesidad de reentrenar o ajustar los parámetros del clasificador cromático para esa región o grupo de dispositivos en particular.
